\newcounter{counterboxspe}
\setcounter{counterboxspe}{0}
%
%\newcounter{counterboxspetosave}
%\setcounter{counterboxspetosave}{0}
%\stepcounter{counterboxspetosave}
%\newcounter{paramFireRegime}
%\setcounter{paramFireRegime}{\thecounterboxspetosave{}}
%\stepcounter{counterboxspetosave}
%\newcounter{ExtremFires}
%\setcounter{ExtremFires}{\thecounterboxspetosave{}}
%\stepcounter{counterboxspetosave}
%\newcounter{pyroCb}
%\setcounter{pyroCb}{\thecounterboxspetosave{}}
%\stepcounter{counterboxspetosave}
%\newcounter{climatechange}
%\setcounter{climatechange}{\thecounterboxspetosave{}}
%\stepcounter{counterboxspetosave}
%\newcounter{scenarClim}
%\setcounter{scenarClim}{\thecounterboxspetosave{}}


\makeatother
\newcounter{counterchapter}
\setcounter{counterchapter}{1}

\chapter{Introduction to the Arctic aerosol and cloud physics}\label{chap1}
\chaptermark{\color{chapcolor} Introduction to the Arctic aerosol and cloud physics}

 \chapquote{"Human beings, who are almost unique in having the ability to learn from the experience of others, are also remarquable for their apparent disinclination to do so."}{Douglas Adams} 
%%%%%%%%%%%%%%%%%%%%%%%%%%%%%%%%%%%%%%%%%%%%%%

\newpage
\vspace*{\fill}
\begin{singlespace}
\textbf{Chapter content}
    
~

\startcontents[chapters]
\printcontents[chapters]{}{1}{}
\end{singlespace} 
\vspace*{\fill}
%%%%%%%%%%%%%%%%%%%%%%%%%%%%%%%%%%%%%%%%%%%%%%
\newpage


%%%%%%%%%%%%%%%%%%%%%%%%%%%%%%%%%%%%%%%%%%%%%%
\section{What does the Earth get from the Sun}
%%%%%%%%%%%%%%%%%%%%%%%%%%%%%%%%%%%%%%%%%%%%%%






%%%%%%%%%%%%%%%%%%%%%%%%%%%%%%%%%%%%%%%%%%%%%%
\section{The Arctic atmosphere}
%%%%%%%%%%%%%%%%%%%%%%%%%%%%%%%%%%%%%%%%%%%%%%

    \subsection{Polar night and stability}


    \subsection{Polar amplification}


    \subsection{Modelling the Arctic climate}

    % question d'ouverture : comment améliorer la modélistation du climat arctique ?

%%%%%%%%%%%%%%%%%%%%%%%%%%%%%%%%%%%%%%%%%%%%%%
\section{Clouds and microphysics}
%%%%%%%%%%%%%%%%%%%%%%%%%%%%%%%%%%%%%%%%%%%%%%

    \subsection{What is a cloud}
    The Earth's atmosphere contains at any moment not less than $~10^{25}$ hydrometeors. The majority are liquid water droplets, and the rest is ice crystals \citep{cantrell_production_2005}.
    This number is to be put into perspective with the overall cloud coverage on the planet, which tends to be of $70 \%$ \cite{}.


    \subsection{Cloud formation and evolution}

from Cantrell and Heymsfield 2005:
In Kurt Vonnegut's novel, Cat's Cradle, Dr.
Asa Breed explains the process using the analogy of
a stack of cannonballs. The cannonballs on the low-
est level "teach" the ones in subsequent layers where
they should be in order to form a stable structure.
Heterogeneous nucleation of ice is envisioned as such
a process, where a substrate forms a template onto
which water molecules adsorb in a configuration
close enough to the crystalline structure of ice that
the energy barrier between phases is substantially
reduced. \footnote{Kurt Vonnegut's brother, Bernard, was one of the first scientists in the United States, along with David Turnbull and Vincent Schaefer, to research the subject of heterogeneous nucleation of ice on a physicochemical level.}

    \subsection{Cloud phase}
    Although the overall cloud ice content is small compared to the liquid content, ice crystals have a great influence on cloud evolution, precipitations and radiative budget.


    \subsection{Ice nucleation}


    % question d'ouverture : comment représenter les INP dans les modèles, et quel est leur impact sur les nuages et le bilan radiatif ?

%%%%%%%%%%%%%%%%%%%%%%%%%%%%%%%%%%%%%%%%%%%%%%
\section{Aerosols particles in the Arctic}
%%%%%%%%%%%%%%%%%%%%%%%%%%%%%%%%%%%%%%%%%%%%%%

    \subsection{Liquid and solid particles in suspension}


    \subsection{Natural or anthropogenic?}


    \subsection{Emission mechanisms}


    % question d'ouverture : quelles sont les sources d'aerosols importantes pour les INPs arctiques ? 

%%%%%%%%%%%%%%%%%%%%%%%%%%%%%%%%%%%%%%%%%%%%%%
\stopcontents[chapters]
